\section{交付与复现}

\subsection{目录结构}

\begin{verbatim}
CSA_HW/
+--  README.md                     # 项目总览
+--  report/                       # 交付层
|   +--  00_选题重定位.md
|   +--  01_报告提纲.md
|   +--  02_结论台账.md
|   +--  03_局限与口径.md
|   +--  04_答辩提纲.md
|   +--  figures/                   # 13 张图（PNG + PDF）
|   \--  methodology/               # 本手册（tex + pdf）
+--  pipeline/                     # 分析管线
|   +--  CONVENTIONS.md             # 口径契约
|   +--  MANIFEST.md                # 45 单元定义
|   +--  csa_common.py              # 共享地基
|   +--  00_build_data.py           # 清洗 + 31 锚点
|   +--  data/                      # 清洗产物
|   \--  out/<单元>/                # result.md + CSV + run_*.py + RUN.txt
+--  verify/                       # 独立校验
+--  USTC网络拓扑分析/              # 原始数据（只读）
\--  archive/                      # 归档（旧预实验、旧进度）
\end{verbatim}

\subsection{运行方式}

\begin{verbatim}
cd pipeline
python 00_build_data.py          # 必须 31/31 PASS
python out/<单元>/run_<xx>.py    # 各单元独立
bash run_all.sh                  # 可选：一键全跑
bash run_all.sh --only E5,F3     # 只跑匹配单元
\end{verbatim}

本机验证环境：Python 3.13.5 / networkx 3.4.2 / pandas 2.2.3 / numpy 2.1.3；
流水线钉死环境：Python 3.12 / networkx 3.6.1 / pandas 3.0.2 / numpy 2.4.4 / scipy 1.17.1 / scikit-learn 1.8.0。
两套环境都通过 31/31 锚点；但 Louvain/PageRank 尾数可能漂，正式引用以仓库版 CSV 为准。

\subsection{锚点与验收}

\begin{itemize}
  \item \code{00\_build\_data.py} 写出 \code{data/anchors.csv}，31 项 PASS。
  \item 每个单元开跑前 \code{assert\_anchors()}；不过就停。
  \item \code{verify/check\_claims.py}：14 项契约断言 + 6 项同学原报对照，退出码 0。
  \item I1--I5：数值一致性、口径违规、统计效度、因果措辞、跨机复现。
\end{itemize}

\subsection{交付物清单}

\begin{table}[H]
\centering
\caption{交付物}
\begin{tabularx}{\textwidth}{@{}p{0.28\textwidth}X@{}}
\toprule
交付物 & 说明 \\
\midrule
报告正文 & \code{report/05} 正文草稿，聚焦结构韧性与失效聚集 \\
结论台账 & \code{report/02}；每条结论的“能说/需限定/不能说” \\
图 & \code{report/figures/} 13 张，含韧性阈值、出链缺失与分层对照；PNG 300dpi + 矢量 PDF \\
代码 & \code{pipeline/}；44 个单元 + 共享地基 + 一键复现 \\
审计 & \code{pipeline/out/I1--I5}；硬伤 0 \\
答辩 & \code{report/04}；13 页提纲 + 30 秒口播 \\
\bottomrule
\end{tabularx}
\end{table}
